\documentclass[a4paper,12pt]{article}

\usepackage[slovene]{babel}
\usepackage[T1]{fontenc}
\usepackage[utf8]{inputenc}
\usepackage{amsmath,amssymb}
\usepackage[margin=2cm,headheight=16pt]{geometry}
\usepackage{enumitem}
\usepackage{xcolor}
\usepackage{fancyhdr}
\usepackage{titlesec}

\definecolor{bordotemno}{RGB}{101,8,34}
\definecolor{bordosvetlo}{RGB}{153,67,88}
\definecolor{svetloroz}{RGB}{255,243,245}

\pagestyle{fancy}
\fancyhf{}
\fancyhead[L]{\small Matematika \textbar\ Funkcije}
\fancyhead[R]{\small Preverjanje znanja}
\fancyfoot[C]{\thepage}
\renewcommand{\headrulewidth}{0.4pt}
\setlist[enumerate,1]{label=\alph*),leftmargin=*,itemsep=5pt,topsep=3pt}
\setlist[itemize]{leftmargin=*,itemsep=2pt,topsep=2pt}
\titleformat{\section}{\large\bfseries\color{bordotemno}}{}{0pt}{}
\titlespacing*{\section}{0pt}{1.3em}{0.5em}

\newcommand{\naloga}[3]{%
	\section*{Naloga #1: #2 \hfill\normalfont\normalsize(#3 točk)}%
}
\newcommand{\odgovor}[1]{\par\noindent\rule{\linewidth}{0.35pt}\par\vspace{#1}}

\begin{document}

\begin{center}
	{\LARGE\bfseries\color{bordotemno} PREVERJANJE ZNANJA}\par
	\vspace{2mm}
	{\large Funkcije \textbar\ Matematika, 2. letnik}
\end{center}

\vspace{3mm}
\noindent
\fcolorbox{bordotemno}{svetloroz}{%
	\parbox{\dimexpr\linewidth-2\fboxsep-2\fboxrule\relax}{%
		\textbf{Ime:} \rule{0.31\linewidth}{0.4pt}\hfill
		\textbf{Priimek:} \rule{0.31\linewidth}{0.4pt}\par\vspace{5mm}
		\textbf{Razred:} \rule{0.20\linewidth}{0.4pt}\hfill
		\textbf{Datum:} \rule{0.25\linewidth}{0.4pt}\hfill
		\textbf{Točke:} \rule{1.1cm}{0.4pt} / 100
	}%
}

\vspace{4mm}
\noindent\textbf{Navodila:} Preberi vsako nalogo in zapiši postopek. Odgovore utemelji, kjer je to zahtevano.\
\textbf{Skupaj: 100 točk.} Nalogi 1 in 2 sta teoretični.

\naloga{1}{}{10}
\begin{enumerate}
	\item Pojasni, kaj mora veljati, da je prirejanje funkcija. \hfill(3 t)
	\item Razloži razliko med definicijskim območjem in zalogo vrednosti funkcije. \hfill(4 t)
	\item Naštej tri načine, na katere lahko podamo funkcijo. \hfill(3 t)
\end{enumerate}
\odgovor{6cm}

\newpage
\naloga{2}{}{10}
\begin{enumerate}
	\item Zapiši pogoj, po katerem prepoznamo sodo oziroma liho funkcijo. \hfill(4 t)
	\item Opiši simetrijo grafa sode in grafa lihe funkcije. \hfill(2 t)
	\item Kaj pomeni, da je funkcija injektivna?  \hfill(2 t)
	\item Kdaj ima funkcija inverzno funkcijo? \hfill(2 t)
\end{enumerate}
\odgovor{7cm}

\newpage
\naloga{3}{}{15}
Naj bo $f(x)=2x-3$ in $g(x)=\sqrt{5-2x}$.
\begin{enumerate}
	\item Izračunaj $f(-1)$ in $f(0)$ ter določi ničlo funkcije $f$. \hfill(6 t)
	\item Določi definicijsko območje funkcije $g$ in njeno ničlo. \hfill(5 t)
	\item Določi definicijsko območje funkcije $r(x)=\dfrac{x+1}{x^2-9}$. \hfill(4 t)
\end{enumerate}
\odgovor{7cm}

\newpage
\naloga{4}{}{20}
Dana je funkcija $p(x)=x^2-4x+3$.
\begin{enumerate}
	\item Določi definicijsko območje in ničle funkcije. \hfill(4 t)
	\item Določi intervale, na katerih je funkcija pozitivna oziroma negativna. \hfill(4 t)
	\item Zapiši funkcijo v temenski obliki ter določi teme in najmanjšo vrednost. \hfill(6 t)
	\item Določi intervala naraščanja in padanja ter zapiši zalogo vrednosti. \hfill(6 t)
\end{enumerate}
\odgovor{7cm}

\newpage
\naloga{5}{}{20}
Funkcijo $h(x)=-2(x-1)^3+4$ dobimo s preoblikovanjem grafa $y=x^3$.
\begin{enumerate}
	\item Opiši zaporedoma vse premike, raztege ali zrcaljenja grafa. \hfill(5 t)
	\item Določi definicijsko območje in zalogo vrednosti funkcije $h$. \hfill(2 t)
	\item Izračunaj začetno vrednost $h(0)$ in določi ničlo funkcije. \hfill(5 t)
	\item V spodnji koordinatni sistem skiciraj graf funkcije $h$. Označi ničlo in presečišče z osjo $y$. \hfill(8 t)
\end{enumerate}
\vspace{2mm}
\begin{center}
\setlength{\unitlength}{1mm}
\begin{picture}(155,67)
	\put(8,20){\vector(1,0){139}}
	\put(64,4){\vector(0,1){58}}
	\multiput(14,18)(10,0){13}{\line(0,1){4}}
	\multiput(62,5)(0,5){12}{\line(1,0){4}}
	\put(146,16){$x$}
	\put(80,62){$y$}
	\put(66,16){$0$}
\end{picture}
\end{center}

\newpage
\naloga{6}{}{25}
\begin{enumerate}
	\item Reši enačbo $\sqrt{2x+3}=x$. Zapiši pogoj za obstoj, preveri dobljene rešitve in navedi končno rešitev. \hfill(10 t)
	\item Dana je funkcija $q(x)=\dfrac{1}{x^2-4}$. Določi njeno definicijsko območje, asimptote in sodost oziroma lihost. Nato reši neenačbo $q(x)>0$. \hfill(10 t)
	\item Ploščina kvadrata je $P$. Zapiši dolžino njegove stranice kot funkcijo ploščine. Nato izračunaj, za koliko se spremeni dolžina stranice, če se ploščina poveča z $9\,\mathrm{m}^2$ na $36\,\mathrm{m}^2$. \hfill(5 t)
\end{enumerate}
\odgovor{7cm}

\vfill
\noindent\textcolor{bordosvetlo}{\rule{\linewidth}{0.6pt}}\\[-1mm]
\small Pred oddajo preveri računske postopke in ali si odgovoril/-a na vse dele nalog.

\end{document}
